\documentclass[twocolumn, 12pt]{article}
\usepackage[utf8]{inputenc}
\usepackage{amsmath}
\usepackage{marvosym}
\usepackage{textcomp}
\usepackage{geometry}
\usepackage{array}
\usepackage{wrapfig}
\usepackage{multirow}
\usepackage{graphicx}
\usepackage{float}
\usepackage{wrapfig}
\usepackage{caption}

\geometry{a4paper, margin=.75in}



\begin{document}
\small{

\section*{Fundamentals of Electronics}
\subsection*{Equations}
     
    \begin{equation*}
        V=\frac{W}{Q} \hspace{15pt} \textbf{F}=\textit{k}_e\frac{Q_1Q_2}{d^2} \hspace{15pt} I=\frac{Q}{t}
    \end{equation*}
    
    \begin{equation*}
        \text{Life of a battery}=\frac{\text{amp-hour rating}}{\text{amps drawn}}
    \end{equation*}
    \begin{equation*}
        R=\rho \frac{l}{A}[1+\alpha_{20}\Delta T]
    \end{equation*}

    \begin{equation*}
        V=IR \hspace{15pt} G=\frac{1}{R} \hspace{15pt} P=\frac{W}{t} \hspace{15pt} 
    \end{equation*}
    
    \begin{equation*}
        P=\frac{W}{t}=\frac{QV}{t}=IV=\frac{V^2}{R}=I^2R
        \end{equation*}
        \begin{equation*}
            W_{\text{in}}=W_{\text{lost}}+W_{\text{stored}}+W_{\text{out}}
        \end{equation*}
        
        \begin{equation*}
            P_{\text{in}}=P_{\text{lost}}+P_{\text{stored}}+P_{\text{out}}
        \end{equation*}
        
        \begin{equation*}
            W_{\text{stored}}=\frac{1}{2}CV^2=\frac{Q^2}{2C} \hspace{15pt} W_{\text{stored}}=\frac{1}{2}LI^2 
        \end{equation*}
        
        \begin{equation*}
            \eta=\frac{W_{\text{out}}}{W_{\text{in}}}=\frac{P_{\text{out}}}{P_{\text{in}}} 
        \end{equation*}
        
        \begin{equation*}
            \eta_{\text{total}}=\eta_1\times \eta_2\times ...\times \eta_n 
        \end{equation*}
        
        \begin{equation*}
            \text{Voltage Regulation}=\frac{V_{\text{no load}}-V_{\text{full load}}} {V_{\text{full load}}}
        \end{equation*}
        
        \begin{equation*}
            \varepsilon =\frac{\textbf{F}}{Q} =\frac{V}{d} \hspace{15pt}  C=\frac{Q}{V}=\varepsilon \frac{A}{d}  \hspace{15pt} \varepsilon_R=\frac{\varepsilon}{\varepsilon_0}
        \end{equation*}
        
        \begin{equation*}
            \tau_C=RC \hspace{15pt} \tau_L=\frac{L}{R} \hspace{15pt} L=\frac{MN^2A}{l}
        \end{equation*}
        
        \begin{equation*}
            i_C=C\frac{d}{dt} v_C \hspace{15pt} v_{L,\text{ avg}}=\frac{\Delta i_L}{\Delta t}
        \end{equation*}
        
        \begin{equation*}
            t=-\tau ln \left(1-\frac{v_C}{V_{\text{max}}}\right) 
        \end{equation*}
        
        \begin{equation*}
            t=-\tau ln \left(1-\frac{i_L}{I_{\text{max}}}\right)
        \end{equation*}
        
        \begin{equation*}
            v_C=v_f-(v_f-v_i)e^{-t/\tau}
        \end{equation*}
        
        \begin{equation*}
            i_L=I_i-(I_f-I_i)e^{-t/\tau}
        \end{equation*}
        
Charging:
        \begin{equation*}
            v_C=E(1-e^{-t/\tau}) \hspace{15pt}i_C=\frac{E}{R_T}e^{-t/\tau}
        \end{equation*}
        
Discharging:
        \begin{equation*}
            v_C=Ee^{-t/\tau}  \hspace{15pt}i_C=\frac{E}{R_T}(1-e^{-t/\tau})
        \end{equation*}
        
        
For Maximum power transfer:
        
        \begin{equation*}
            R_L=R_{TH} \hspace{15pt} P_{L,\text{ max}} =\frac{V_{TH}^2}{4R_{TH}}
        \end{equation*}
        
\subsubsection*{Alternating Current}
        \begin{equation*}
            I_{\text{rms}} = \frac{I_{\text{max}}}{\sqrt{2}} \hspace{15pt}  V_{\text{rms}} =\frac{V_{\text{max}}}{\sqrt{2}} 
        \end{equation*}
        
        \begin{equation*}
            \omega=2\pi f  \hspace{15pt} T=\frac{1}{f} \hspace{15pt} \textbf{Y}=\frac{1}{\textbf{Z}}  
        \end{equation*}
        
        
        \begin{equation*}
            X_L= \omega L \hspace{15pt} X_C=\frac{1}{\omega C} \hspace{15pt}B=\frac{1}{X}
        \end{equation*}
    
        \(G\) = (\text{sum of areas})/(\text{length of the curve})
        
        \begin{equation*}
            pf=cos|\theta_v-\theta_i| \hspace{15pt} P_{\text{avg}}=V_{\text{rms}}I_{\text{rms}}pf
        \end{equation*}
        
        \begin{center}
            pf leading: the current leads the voltage\\
            pf lagging: the current lags the voltage
        \end{center}
        

\subsubsection*{Impedance}
        The complex opposition to changing current/voltage.
    
    \begin{equation*}
        \textbf{Z}_C=R_C+jX_C \hspace{15pt} \textbf{Z}_L=R_L+jX_L
    \end{equation*}
        \noindent
        Where:
        \begin{itemize}
            \item R is the real form (resistance) [\(\Omega\)]
            \item X is the imaginary form (reactance) [\(\Omega\)]
        \end{itemize}
       
        
\subsubsection*{Admittance}
        The complex allowance of changing current/voltage.
        \begin{equation*}
            \textbf{Y}_C= G_C+jB_C \hspace{15pt} \textbf{Y}_L = G_L +jB_L
        \end{equation*}
        \noindent
        Where:
        \begin{itemize}
            \item G is the real form (conductance) [S]
            \item B is the imaginary form (susceptance) [S]
        \end{itemize}
        
        
        
\subsubsection*{Series Circuits} 

         \begin{equation*}
            \sum V_{\text{closed loop}}=0 \hspace{15pt} \frac{V_x}{V_S}=\frac{R_x}{R_T}
        \end{equation*}
        
        \begin{equation*}
            R_{T}=R_1+R_2+...+R_n \hspace{15pt} 
        \end{equation*}
        
        \begin{equation*}
            \frac{1}{C_T}=\frac{1}{C_1}+\frac{1}{C_2}+...+\frac{1}{C_N}
        \end{equation*}
        
        \begin{equation*}
            L_T=L_1+L_2+...+L_N
        \end{equation*}
        
        \begin{equation*}
            \textbf{Z}_T = \textbf{Z}_1 + \textbf{Z}_2 +...+\textbf{Z}_N
        \end{equation*}
        
        
\subsubsection*{Parallel Circuits}
        
        \begin{equation*}
            \sum I_{\text{in}}= \sum I_{\text{out}} \hspace{15pt}  I_S=I_1+I_2+...+I_n               
        \end{equation*}
        
        \begin{equation*}
            I_S R_T=I_x R_x \hspace{10pt}   R_T=\frac{R_1R_2}{R_1+R_2} 
        \end{equation*}
        
        \begin{equation*}
            P_S=P_{R_1}+P_{R_2}+...+P_{R_n}
        \end{equation*}
        
        \begin{equation*}
            \frac{1}{R_T}=\frac{1}{R_1}+\frac{1}{R_2}+...+\frac{1}{R_n}
        \end{equation*}
        
        \begin{equation*}
            C_T=C_1+C_2+...+C_N
        \end{equation*}
        
        \begin{equation*}
            \frac{1}{L_T}=\frac{1}{L_1}+\frac{1}{L_2}+...+\frac{1}{L_N}
        \end{equation*}
        
        \begin{equation*}
            \frac{1}{\textbf{Z}_T} = \frac{1}{\textbf{Z}_1}+ \frac{1}{\textbf{Z}_2}  +...+\frac{1}{\textbf{Z}_N} 
        \end{equation*}
        
        
%\subsection*{Source Conversion}
  

\subsection*{Network Theorems}

    \subsubsection*{Superposition}
    \begin{enumerate}
        \item Activate one source at a time
        \begin{enumerate}
            \item Voltage source: short it
            \item Current source: open it
        \end{enumerate}
        \item Find \(I'_1\), \(I'_2\), etc.
        \item Add I primes together to find current values
    \end{enumerate}
    
    \subsubsection*{Mesh Current Analysis}
    \begin{enumerate}
        \item Assume current directions
        \item Write an equation per loop using KVL
        \item Solve for unknown current(s)/value(s)
    \end{enumerate}
    
    \subsubsection*{Super Mesh}
        \begin{center}
                 *Only if you have a current source in the middle*
        \end{center}
        \begin{enumerate}
            \item Asssume current directions 
            \item Deactivate (open) the current source
            \item Write one equation using both \(I_1\) and \(I_2\)
            \item For the other equation, use the sum of \(I_1\) and \(I_2\)
            \item Solve the system of equations
        \end{enumerate}
        
    \subsubsection*{Nodal Analysis}
    \begin{enumerate}
        \item Determine the number of nodes 
        \item Pick and label each node
        \item Apply KCL at each node assuming each node has largest voltage value
        \item Solve system of equations
        
    \end{enumerate}
    
    \subsubsection*{Super Node}
        \begin{center}
             *only if you have a voltage source preventing nodal analysis*
        \end{center}
        \begin{enumerate}
            \item Short the source
            \item Proceed as though nodal analysis would work
            \item Solve the system of equations
        \end{enumerate}
    \subsubsection*{Thevenin's Theorem}
        \begin{enumerate}
            \item Identify and remove the load
            \item Find \(\ R_{TH}\) by deactivating all sources simultaneously
            \item Find \(\ V_{TH}\) across terminals from step 1
            \item Draw the equivalent circuit
        \end{enumerate}  
        
    \subsubsection*{Norton's Theorem}    
    \begin{enumerate}
        \item Identify and remove the load
        \item Find \(\ R_N\) by deactivating all sources simultaneously
        \item Find\(\ V_N \) by shorting terminals from step 1
        \item Draw the equivalent Norton circuit 
    \end{enumerate}
        
\end{document}